\chapter{TỔNG QUAN BÀI TOÁN VÀ NHIỆM VỤ CỦA TV2}
\label{chap:tongquan}

\section{Bài toán định tuyến xe xanh đa mục tiêu có khung thời gian}

Một công ty có một kho hàng và một đội xe tải giống nhau, cần giao hàng cho $n$ khách. Mỗi khách $i$
có toạ độ, lượng hàng cần nhận $q_i$, thời gian dỡ hàng $s_i$ và hai khung thời gian lồng nhau: khung kỳ
vọng $[ETa_i, LTb_i]$ là khoảng khách muốn nhận hàng nhất, khung chấp nhận được $[ET_i, LT_i]$ rộng hơn,
ngoài khung này khách không hài lòng chút nào. Mỗi xe xuất phát từ kho, đi qua một số khách theo một thứ
tự rồi quay về kho. Tổng lượng hàng trên một xe không vượt quá tải trọng $Q$, và số xe dùng không vượt
quá $K$. Lời giải là cách chia khách cho các xe và thứ tự ghé thăm trên từng xe.

Bài toán có ba mục tiêu cùng lúc:
\begin{itemize}
  \item $f_1$: tổng chi phí giao hàng (CNY), cần cực tiểu;
  \item $f_2$: tổng lượng khí thải CO\textsubscript{2} (kg), cần cực tiểu;
  \item $f_3$: mức hài lòng trung bình của khách, trong đoạn $[0,1]$, cần cực đại.
\end{itemize}
Ba mục tiêu này xung đột nhau. Dùng thêm xe giúp đến đúng giờ hơn (tăng $f_3$, giảm phạt trễ trong $f_1$)
nhưng làm tăng chi phí cố định và quãng đường (tăng $f_2$). Vì vậy kết quả không phải một lời giải duy
nhất mà là một tập lời giải không trội nhau, gọi là tập Pareto.

\subsection{Ba hàm mục tiêu}

Các tham số được lấy theo bảng 2 của bài báo gốc và ghi trong \code{config.yaml}.

\paragraph{Chi phí $f_1$.} Gồm bốn thành phần:
\begin{equation}
  f_1 = F_1 + F_2 + F_3 + F_4 .
\end{equation}
$F_1 = \sigma_1 \cdot \text{Fuel}$ là tiền nhiên liệu, với $\sigma_1 = 10$ CNY/lít. Suất tiêu thụ trên
một chặng phụ thuộc tải hiện có trên xe trước khi giao cho khách tiếp theo:
\begin{equation}
  \rho(L) = \rho_0 + \frac{\rho^* - \rho_0}{Q}\, L ,
  \qquad \rho_0 = 0{,}122 \text{ L/km},\; \rho^* = 0{,}388 \text{ L/km},
\end{equation}
nên lượng nhiên liệu của chặng dài $d$ là $d\cdot\rho(L)$. $F_2 = \sigma_2 m$ là chi phí cố định cho
$m$ xe được dùng ($\sigma_2 = 100$ CNY/xe). $F_3 = \sigma_3 \sum d$ là chi phí theo quãng đường
($\sigma_3 = 2$ CNY/km). $F_4$ là tiền phạt lệch khung kỳ vọng:
\begin{equation}
  F_4 = \frac{\alpha \sum_i \max(ETa_i - t_i, 0) + \beta \sum_i \max(t_i - LTb_i, 0)}{60},
\end{equation}
với $t_i$ là giờ xe đến khách $i$ tính bằng phút, $\alpha = 20$ CNY/giờ cho đến sớm, $\beta = 60$
CNY/giờ cho đến muộn. Phép chia cho 60 đổi số phút sang giờ.

\paragraph{Khí thải $f_2$.} Theo mô hình MEET, với vận tốc không đổi $v = 40$ km/h:
\begin{equation}
  f_2 = \frac{\varphi(v)\,\psi(v)}{1000} \sum d, \qquad
  \varphi(v) = c_0 + c_1 v^3 + \frac{c_2}{v}, \qquad
  \psi(v) = d_0 + d_1 v + d_2 v^2 + \frac{d_3}{v}.
\end{equation}
Bài báo in số hạng $d_2 v^3$ khiến $\psi$ âm ở $v = 40$; nhóm sửa thành $d_2 v^2$, khi đó mỗi km thải
khoảng $679{,}44$ g CO\textsubscript{2}. Hệ quả quan trọng là $f_2$ tỉ lệ thuận với tổng quãng đường,
nên $f_2$ và thành phần $F_3$ của $f_1$ biến thiên cùng chiều.

\paragraph{Mức hài lòng $f_3$.} Mỗi khách có hàm hài lòng hình thang $W_i(t)$, thang 0--100:
\begin{equation}
  W_i(t) =
  \begin{cases}
    0 & t < ET_i \text{ hoặc } t > LT_i,\\[2pt]
    100\,\dfrac{t - ET_i}{ETa_i - ET_i} & ET_i \le t < ETa_i,\\[6pt]
    100 & ETa_i \le t \le LTb_i,\\[2pt]
    100\,\dfrac{LT_i - t}{LT_i - LTb_i} & LTb_i < t \le LT_i,
  \end{cases}
  \qquad
  f_3 = \frac{1}{n}\sum_{i=1}^{n} \frac{W_i(t_i)}{100}.
\end{equation}
Nhóm tính $W$ tại giờ xe đến (\code{satisfaction\_time: arrival}). Xe đến sớm thì chờ tới $ETa_i$ rồi
mới dỡ hàng.

\subsection{Dữ liệu}

Nhóm dùng bộ Solomon~\cite{solomon1987} (bản 100 khách, lấy $n$ khách đầu cho các cỡ 25 và 50), tải trọng $Q = 400$ kg,
đội xe tối đa $K = 25$, khung chấp nhận được là khung gốc nới thêm 30 phút mỗi đầu. Ngoài ra có bộ dữ
liệu thực tế 29 khách (\code{case29}) từ bảng 8 của bài báo, tải trọng 4000 kg, chưa giới hạn số xe.
Phần việc của TV2 thực nghiệm trên bốn instance đại diện: C102-25 (khách phân cụm, khung hẹp), R201-50
(phân bố ngẫu nhiên, khung rộng), RC201-100 (hỗn hợp, cỡ lớn) và case29-29 (dữ liệu thực tế).

\section{Mã hoá lời giải và bộ giải mã}
\label{sec:decoder}

Một cá thể là một hoán vị $\pi$ của các khách $1,\dots,n$, không chứa kho. Bộ giải mã (\code{decoder.py})
cắt hoán vị thành các tuyến theo đúng thứ tự trong hoán vị. Nhóm dùng quy tắc cắt
\code{capacity\_time}: đi dọc hoán vị, mở tuyến mới khi khách tiếp theo làm xe vượt tải, hoặc khi xe sẽ
đến khách tiếp theo sau $LT$ mà đội xe vẫn còn xe trống. Sau khi cắt, bộ giải mã đi dọc từng tuyến, tính
giờ đến, tải, quãng đường rồi ra $(f_1, f_2, f_3)$ cùng các thông tin phụ như danh sách tuyến, số xe và
cờ \code{fleet\_ok} cho biết số xe có vượt $K$ hay không.

Cách mã hoá này luôn cho lời giải hợp lệ về tải, nhưng có một hệ quả mà phần tìm kiếm cục bộ phải xử lý:
chỗ cắt tuyến phụ thuộc vào toàn bộ thứ tự trong hoán vị. Đổi thứ tự khách bên trong một tuyến có thể
làm chỗ cắt dịch đi, tức bộ giải mã tạo ra một bộ tuyến khác (mục~\ref{sec:drift}).

\section{Thuật toán NSGA-II của nhóm}

TV1 cài đặt NSGA-II gốc trong \code{src/nsga2.py}. Quần thể có $N_{pop} = 80$ cá thể, chạy $G = 200$
thế hệ. Mỗi thế hệ: chọn cha mẹ bằng giải đấu nhị phân theo hạng và khoảng cách đông đúc, lai ghép thứ
tự (OX) với xác suất $p_c = 0{,}7$, đột biến đảo đoạn hoặc hoán đổi hai vị trí với xác suất
$p_m = 0{,}1$ (Solomon) hoặc $0{,}2$ (case29), đánh giá con, gộp cha mẹ với con rồi chọn lọc môi trường
theo front và khoảng cách đông đúc, có loại các hoán vị trùng nhau. Ràng buộc số xe được xử lý bằng trội
có ràng buộc: lời giải vượt ít xe hơn luôn thắng. Toàn bộ thuật toán dùng một bộ sinh số ngẫu nhiên
\code{np.random.default\_rng(seed)} duy nhất, nên cùng seed luôn cho cùng kết quả.

\section{Bốn cải tiến và nhiệm vụ của TV2}

\begin{table}[H]
\centering
\caption{Bốn cải tiến của nhóm và người phụ trách}
\label{tab:caitien}
\begin{tabularx}{\textwidth}{@{}lXl@{}}
\toprule
\textbf{Cải tiến} & \textbf{Ý tưởng} & \textbf{Người làm} \\
\midrule
\code{init\_smart}   & khởi tạo quần thể bằng láng giềng gần nhất, sắp theo khung giờ và ngẫu nhiên & TV2 \\
\code{local\_search} & 2-opt và Or-opt trên từng tuyến sau đột biến & TV2 \\
\code{diversity}     & SDE hoặc mật độ kNN thay khoảng cách đông đúc, thêm kho lưu trữ ngoài & TV3 \\
\code{dp\_split}, \code{adaptive\_rates} & cắt tuyến tối ưu bằng quy hoạch động; $p_c$, $p_m$ thích nghi & TV4 \\
\bottomrule
\end{tabularx}
\end{table}

Theo phân công trong README của nhóm, TV2 phải làm các việc sau.
\begin{enumerate}
  \item \textbf{Giai đoạn 2 (12/11--18/11).} Cài đặt \code{init\_smart} và \code{local\_search}, mỗi cải
        tiến một tệp trong \code{src/improvements/}, bật tắt bằng một cờ trong \code{config.yaml}, mặc định
        tắt. Viết test riêng. Chứng minh khi tắt cờ thì ra đúng kết quả của thuật toán gốc với cùng seed.
        Chạy 10 seed trên ít nhất 3 instance và so HV với thuật toán gốc.
  \item \textbf{Giai đoạn 3 (19/11--25/11).} Hỗ trợ ghép các nhánh, chuẩn bị mốc hồi quy: tắt hết cờ phải
        ra đúng từng con số như thuật toán gốc.
  \item \textbf{Giai đoạn 4 (26/11--02/12).} Vẽ biểu đồ: front Pareto ba chiều, đường hội tụ HV, biểu đồ
        hộp HV và IGD qua 10 seed.
\end{enumerate}

Ngoài ra có bốn quy tắc chung bắt buộc cho mọi cải tiến: (i) cờ tắt thì code cải tiến không được gọi
bộ sinh số ngẫu nhiên lần nào, (ii) chỉ dùng bộ sinh do \code{nsga2.run} truyền vào, (iii) mọi tham số
đặt trong \code{config.yaml} dưới mục \code{improvements}, (iv) mọi lần đánh giá thêm phải được cộng vào
số lần đánh giá \code{n\_evals}, để so sánh giữa các cấu hình được công bằng.

Riêng với tìm kiếm cục bộ, README nêu sẵn một điểm phải cẩn thận: sau khi cải thiện một tuyến rồi nối
các tuyến lại thành hoán vị, bộ giải mã cắt lại có thể ra bộ tuyến khác, khiến phần cải thiện bị mất
hoặc bị tính sai. Chương~\ref{chap:ls} trình bày cách xử lý điểm này.
